%%%%%%%%%%  scan idx 680  |  folio 667  %%%%%%%%%%
\SGAPrintPage{667}%
\sgaStableTarget{sga6.e14.expose}\section*{\underline{Exposé XIV}\\[3pt] \underline{\underline{OPEN PROBLEMS IN INTERSECTION THEORY}}}
\begin{center}
\underline{by A. Grothendieck}
\end{center}

In this exposé we review some of the problems connected with the questions encountered in this seminar that deserve to be studied. We leave aside the problems connected with Serre's \uline{local theory} of intersections [14], which we have not had occasion to use in the Seminar, and refer in particular to loc. cit. for the two fundamental conjectures that remain open in this direction (loc. cit. Chap. V \S~4). We also pass over the problems, very important for arithmetic applications, connected with the \uline{structure of groups of cycle classes}, especially on a smooth projective algebraic scheme; for a discussion of these problems we refer to S. KLEIMAN, Algebraic cycles and the Weil conjectures in Dix exposés sur la cohomologie des schémas (Masson-North Holland Publishing Cie). In no.~4 of this exposé we also give some supplements to the Seminar, concerning in particular the \uline{Chow ring}, which had not found a place in the preceding exposés.

\sgaStableTarget{sga6.e14.sec.1}\section*{1. Operations $\Lambda^i$ in the derived category $D(X)$\SGAAuthorFootMark{*}}
\SGAAuthorFootText{*}{cf. the note on page 23.}

Let $X$ be a locally ringed topos, and let $K^\bullet(X)$ be the ring of classes of perfect complexes on $X$ (Exp. IV). We would like to define on $K^\bullet(X)$ a $\lambda$-ring structure (Exp. V), functorial in $X$, such that the natural homomorphism
\[
 K^\bullet_{\mathrm{naif}}(X)\longrightarrow K^\bullet(X)
\]
(where the first member is the ring of classes of locally free Modules on $X$, endowed with the $\lambda$-structure defined in Exp. V) be a $\lambda$-homomorphism. For this purpose, one would need to define functors (generally nonadditive) of ``exterior power''
\[
 \Lambda^i:D^-(X)\longrightarrow D^-(X),
\]
taking pseudo-coherent (resp. perfect) complexes to pseudo-coherent (resp. perfect) complexes and having the formal properties needed, with respect to exact triangles, to define a $\lambda$-structure on $K^\bullet(X)$ by the formula
\sgaStableTarget{sga6.e14.eq.1-1}\begin{equation}\tag{1.1}
 \lambda^i(\mathrm{c}\ell^\bullet(L_\bullet))=\mathrm{c}\ell^\bullet(\Lambda^i L_\bullet).
\end{equation}

%%%%%%%%%%  scan idx 681  |  folio 668  %%%%%%%%%%
\SGAPrintPage{668}%
When $X$ is of characteristic zero, i.e. its structure sheaf is a sheaf of $Q$-algebras, the problem is solved simply by taking
\sgaStableTarget{sga6.e14.eq.1-2}\begin{equation}\tag{1.2}
 \Lambda^i L_\bullet=\bigl(\overset{i}{\otimes}{}^{\mathbb L}L_\bullet\bigr)^{\mathrm{alt}},
\end{equation}
where the $i$-th tensor power (in the sense of derived categories) of $L_\bullet$ is regarded as an object of $D^-(\underline{O}_i)$, with $\underline{O}_i$ the group algebra of the symmetric group $\mathfrak S_i$ with coefficients in $\underline{O}_X$, and where the superscript alt denotes passage to the ``alternating part under the action of $\mathfrak S_i$'' (an operation that makes sense by the characteristic-zero hypothesis). A formula of this type was already used in [8] (cited below as [RRR]) for the first proof of the Riemann--Roch theorem.

Without a hypothesis on the characteristic, when $L_\bullet$ is zero in degrees $>0$, Dold--Puppe [6] shows that $L_\bullet$ is canonically homotopic to the complex associated with a semi-simplicial Module $L'_\bullet$ canonically associated with $L_\bullet$, and a fairly natural candidate for $\Lambda^iL_\bullet$ (for $L_\bullet$ flat, say) would be the complex associated with the semi-simplicial Module $\Lambda^iL'_\bullet$ (whose component of degree $n$ is $\Lambda^i L^{\prime}_n$). But to extend the definition to complexes of $D^-(X)$ no longer assumed to be zero in degrees $>0$, one would have to study the behaviour of the preceding operation with respect to the translation operation $L_\bullet\mapsto L_\bullet[1]$ (or \uline{suspension}, in the terminology of [6]); unfortunately, there does not seem to be an isomorphism $\Lambda^i(L_\bullet[1])\simeq(\Lambda^i L_\bullet)[i]$ that would make the desired extension trivial.

\sgaStableTarget{sga6.e14.sec.2}\section*{2. The Riemann--Roch formula without projectivity hypotheses}

\sgaStableTarget{sga6.e14.para.2-1}\textbf{2.1.} Once a natural $\lambda$-structure on the $K^\bullet(X)$ is available, it is possible to formulate the Riemann--Roch theorem, in the form considered in Exp. VIII, for any proper complete-intersection morphism $f:X\to Y$ (and even, more generally, for a relative scheme $X$ over a locally ringed topos $Y$, with $X$ still subject to the corresponding relative conditions). The Todd class of $f$ is defined by means of the class in $K^\bullet(X)$ of the relative cotangent complex $T_f$ for $f$, constructed in Exp. 0 4.4; this complex is perfect here by the local hypothesis on $f$. It seems clear that a proof of the Riemann--Roch formula in this general case will require the introduction of essentially new ideas, and it is possible that these same ideas could provide the key
%%%%%%%%%%  scan idx 682  |  folio 669  %%%%%%%%%%
\SGAPrintPage{669}%
to a proof of the analogous formula in the complex analytic (or rigid analytic!) case. Recall that in the complex analytic case, only the case of a locally free sheaf $F$ on a compact nonsingular variety $X$ over a point $Y$ is currently known (by the Atiyah--Singer index formula); the case in which $F$ is merely assumed coherent still seems to lie beyond the techniques known so far.

\sgaStableTarget{sga6.e14.para.2-2}\textbf{2.2.} As an essential ingredient of the Riemann--Roch formula, let us also recall the \uline{finiteness problem} raised in Exp. III: for a proper morphism $f:X\to Y$ and a complex $L_\bullet$ on $X$ that is pseudo-coherent relative to $Y$, is $Rf_*(L_\bullet)$ pseudo-coherent? A general proof of this conjecture will probably use algebraic ideas that will make it possible to prove Grauert's finiteness ``theorem'' [7] under considerably more general conditions than those of loc. cit., starting directly with topological $\C$-algebra structures on base spaces (or topoi) $Y$ more general than analytic spaces (including the case of a locally compact space $Y$ with the sheaf of continuous complex functions, or a $C^m$ manifold with the sheaf of $C^m$ complex functions, or a real analytic space, etc.).

\sgaStableTarget{sga6.e14.para.2-3}\textbf{2.3.} Finally, in these questions one would need to replace the properness hypothesis on the morphism $f$ by the hypothesis that the cohomology sheaves $H^i(L_\bullet)$ of the complexes under consideration on $X$ have their supports proper over $Y$.

\sgaStableTarget{sga6.e14.sec.3}\section*{3. A Riemann--Roch formula ``without denominators'' for an immersion}

\sgaStableTarget{sga6.e14.para.3-1}\textbf{3.1.} Let $f:Y\to X$ be a regular closed immersion (Exp. VII), with conormal sheaf $\underline N$. For simplicity, suppose that $X$ admits an ample invertible sheaf. The Riemann--Roch formula for $y\in K^\bullet(Y)$:
\[
 \ch_X(f_*(y))=f_*\bigl(\ch_Y(y)\operatorname{Todd}(-\check{\underline N})\bigr),
\]
together with the general formula (Exp. VII):
\[
 f_*(y')\,f_*(y'')=f_*\bigl(y'y''c_d(\underline N)\bigr),\qquad y',y''\in \operatorname{Gr}(Y)_{\underline Q},
\]
(with $f$ assumed to have codimension $d$), makes it possible to express the Chern classes of $f_*(y)$ in $\operatorname{Gr}^\bullet(X)_\Q$ in the form
%%%%%%%%%%  scan idx 683  |  folio 670  %%%%%%%%%%
\SGAPrintPage{670}%
\sgaStableTarget{sga6.e14.eq.3-1}\begin{equation}\tag{3.1}
 c_i(f_*(y))=f_*\bigl(P_i((c_\alpha(y))_\alpha,(c_\beta(\underline N))_\beta)\bigr),
\end{equation}
where $P_i$ is a universal polynomial in the $c_\alpha(y)$, $c_\beta(\underline N)$, a polynomial whose coefficients are a priori rational. In fact, one finds [RRR] that these polynomials have \uline{integer} coefficients. This raises the question of the validity of a Riemann--Roch formula ``without denominators''. In fact, such a formula for $c_i(f_*(y))$ cannot in general be sought in $\operatorname{Gr}^\bullet(X)$ itself, since the homomorphism $f_*:K^\bullet(Y)\to K^\bullet(X)$ is not itself compatible with the $\lambda$-filtrations up to a shift by $d$ (4.6), but is so only up to torsion; thus one cannot define a homomorphism $f_*:\operatorname{Gr}^\bullet(Y)\longrightarrow \operatorname{Gr}^\bullet(X)$ giving the natural homomorphism $f_*:\operatorname{Gr}^\bullet(Y)_\Q\to\operatorname{Gr}^\bullet(X)_\Q$ and allowing formula (3.1) to be interpreted as a formula in $\operatorname{Gr}^\bullet(X)$, and not merely in $\operatorname{Gr}^\bullet(X)_\Q$. But when $X$ is assumed regular, so that $K^\bullet(X)\xrightarrow{\sim} K_\bullet(X)$ (IV 2.5), one can endow $K^\bullet(X)$ with the \uline{topological filtration} by the codimension of supports, inherited from the corresponding filtration of $K_\bullet(X)$. It is plausible that this filtration is compatible with the ring structure, a fact that can be proved at least [RRR] when $X$ is smooth over a field $k$, using the theory of the Chow ring [4]. When this is so, the associated graded
\[
 \Gtop^\bullet(X)
\]
is a \uline{graded ring}, and, also assuming $Y$ regular, we have done what is needed to obtain a natural homomorphism of degree $d$ of graded groups
\[
 f_*=f_*^{\mathrm{top}}:\Gtop^\bullet(Y)\longrightarrow \Gtop^\bullet(X).
\]
Since the $\lambda$-filtration of $K^\bullet(X)$ is finer than the topological filtration by codimension just considered (X 1.3.3), one obtains a canonical homomorphism of graded groups (and even of graded rings)
\[
 \operatorname{Gr}^\bullet(X)\longrightarrow \Gtop^\bullet(X),
\]
which is moreover an isomorphism mod torsion (Exp. VIII). This allows us, as in [RRR], to consider a theory of Chern classes on $X$ with values in $\Gtop^\bullet(X)$, and likewise on $Y$, and thus to interpret formula (3.1) as a formula in $\Gtop^\bullet(X)$. In this form, the formula is proved in [RRR] when $X$ is of \uline{characteristic zero}, at least when $X$ is smooth over a field (a hypothesis serving only to ensure that the topological filtration of $K^\bullet(X)$ is compatible with the ring structure). The characteristic-zero hypothesis is used essentially
%%%%%%%%%%  scan idx 684  |  folio 671  %%%%%%%%%%
\SGAPrintPage{671}%
to obtain formulas of the type (1.1), (1.2), and essentially the same proof should apply without a hypothesis on the characteristic, once the question raised in no.~1 has been answered affirmatively.

\sgaStableTarget{sga6.e14.para.3-2}\textbf{3.2.} When $X$ is smooth over a field, one also has the Chow ring
\[
 A(X),
\]
and a theory of Chern classes on $X$ with values in $A(X)$ [9], finer than the theory with values in $\Gtop^\bullet(X)$ just introduced---these latter Chern classes being the images of the Chow-style Chern classes under the canonical surjective homomorphism (cf. 4.1)
\[
 \phi:A(X)\longrightarrow \Gtop^\bullet(X).
\]
This homomorphism is an isomorphism mod. torsion, but not an isomorphism in general (4.7). With this understood, for smooth $X$ and $Y$ the inclusion $f:Y\to X$ again defines a homomorphism of degree $d$ of graded groups
\[
 f_*:A(Y)\longrightarrow A(X),
\]
which allows formula (3.1) to be interpreted as a formula in $A^i(X)$\SGAAuthorFootnote{*}{The formula, as well as the variant mentioned in 3.3, has since been proved by P. JOUANOLOU (Nov.\ 1969). Cf.}. This formula reduces to $0=0$ for $i<d$ and is therefore of interest only for $i\geq d$. Even for $i=d=2$ this formula has not been proved at present (cf. 4.3), nor has the less precise formula deduced from it by replacing rational equivalence of algebraic cycles with algebraic equivalence. Unlike the analogous formula in $\Gtop^\bullet(X)$, there seems to be no very convincing reason to suppose that the formula without denominators must hold in $A(X)$. In view of our almost total ignorance of the nature of the groups $A^i(X)$ for $i\geq 2$, a good way to test it seems to be to test the formulas that can be deduced from it in groups $A^1(Z)=\Pic(Z)$, for example by applying $g_*:A(X)\longrightarrow A(Z)$ to both members of (3.1), where $g:X\to Z$ is a proper morphism into a smooth quasi-projective $Z$ of dimension $\dim X-i+1$.

\sgaStableTarget{sga6.e14.para.3-3}\textbf{3.3.} Finally, to give a meaning to (3.1), one can also use the cohomological theory of Chern classes (SGA 5 VII), whenever \uline{Gysin homomorphisms}
\[
 f_*:H^p(Y,\mu_n^{\otimes i})\longrightarrow H^{p+2d}(X,\mu_n^{\otimes(i+d)})
\]
are available (cf. no.~6); this is the case at least when $Y$ and $X$ are both smooth over a scheme $S$ (by the theorem of relative cohomological purity
